\section*{Tóm tắt}
\addcontentsline{toc}{section}{Tóm tắt}

Việc xây dựng ngữ liệu song song ở mức câu giữa văn bản Hán Nôm và bản phiên âm Quốc ngữ từ các ấn bản in khắc đầu thế kỷ XX đòi hỏi sự phối hợp giữa xử lý ảnh tài liệu, nhận dạng quang học và đối chiếu ngôn ngữ lịch sử. Khi làm việc trực tiếp với ảnh quét nguyên bản, quy trình tự động thường gặp nhiều trở ngại đan xen như sai lệch thứ tự trang do đóng sách hoặc quét thiếu, bố cục đọc ngược chiều giữa hai hệ chữ viết, sự thiếu nhất quán về dấu câu trong cổ văn, hiện tượng đa âm cùng dị thể chữ Nôm, cũng như chính tả Quốc ngữ cổ có nhiều hình thái khác với tiếng Việt hiện đại. Trong báo cáo này, nhóm trình bày phương pháp và chuỗi thuật toán áp dụng trên bộ tác phẩm Công giáo song ngữ gồm hai tập (\textit{Sách Dạy Đấng Làm Người Sửa Mình}, Tập 1 và Tập 2, tổng cộng $1.321$ trang quét nguyên bản).

Cách tiếp cận của nghiên cứu được tổ chức thành một quy trình liên kết theo năm bước xử lý chính. Trước hết, cấu trúc vật lý của hai tập sách được phục hồi thông qua mô hình ánh xạ trang tuyến tính từng đoạn có hiệu chỉnh dịch pha và hoán vị tay sách, kết hợp với bước lọc nhiễu viền dựa trên thống kê trung vị để tách phần chính văn. Tiếp theo, không gian tự hình Hán Nôm được chuẩn hóa bằng cách truy nguyên các mã vùng sử dụng riêng (PUA) sang chuẩn Unicode mở rộng theo cơ sở dữ liệu Unihan và chuỗi mô tả cấu trúc biểu ý (IDS) \cite{unihan2025,morioka2008ids}, đồng thời điều chỉnh $19$ quy luật nhận dạng nhầm nét lặp lại trong sách. Ở vế Quốc ngữ, nhóm kết hợp mô hình ngôn ngữ -- thị giác (VLM) giải mã từng dòng nội suy với cơ chế hợp nhất giả thuyết ROVER có ràng buộc ngữ âm lấy từ trang Hán Nôm đối ứng \cite{fiscus1997rover,bai2025qwenvl}, qua đó giảm $77{,}1\%$ số âm tiết nhận dạng sai so với hệ thống truyền thống mà vẫn bảo lưu lớp chính tả cổ. Sau khi chuẩn hóa văn bản ở cả hai vế, bài toán dóng hàng được giải bằng quy hoạch động phân cấp hai tầng: tầng đầu tiên chiếu ranh giới câu từ Quốc ngữ sang dòng Hán Nôm liên tục ở cấp cụm trang nhằm hạn chế trượt lệch tích lũy, còn tầng thứ hai đối chiếu ở mức tự hình -- âm tiết theo khoảng cách Levenshtein có trọng số ngữ âm \cite{levenshtein1966}. Cuối cùng, đối với hai trang bị khuyết bản Quốc ngữ khi số hóa ở Tập 1, hệ thống ước lượng bản phiên âm dự đoán trực tiếp từ phân phối đồng xuất hiện chữ -- âm trong chính cuốn sách. Đánh giá trên toàn bộ $1.319$ trang hiệu lực ($636$ đơn vị ánh xạ chính văn, $5.096$ cặp câu song ngữ 1--1 và $146.277$ tự hình Hán Nôm) ghi nhận độ tương đồng ngữ âm chính thức đạt $98{,}14\%$ trên Tập 1 và $98{,}45\%$ trên Tập 2, với tỷ lệ hàng không so sánh bằng $0{,}00\%$ trên cả hai tập, đều cao hơn ngưỡng yêu cầu $95\%$.

\section{Đặt Vấn Đề và Đặc Trưng Ngôn Ngữ -- Văn Bản Học}

\subsection{Bản chất bài toán và đối tượng nghiên cứu}

Ngữ liệu song song Hán Nôm -- Quốc ngữ là nguồn dữ liệu nền tảng phục vụ các tác vụ phiên âm tự động, nghiên cứu lịch sử tiếng Việt, tra cứu từ nguyên và số hóa di sản thành văn. Đối tượng khảo sát trong báo cáo này là hai tập ấn bản song ngữ đầu thế kỷ XX thuộc thể loại giáo lý -- tu đức Công giáo mang tên \textit{Sách Dạy Đấng Làm Người Sửa Mình}, trong đó Tập 1 gồm $636$ trang quét nguyên bản và Tập 2 gồm $685$ trang quét nguyên bản. Xét theo tiêu chí dịch thuật, toàn bộ nội dung chính văn của cả hai tập đều thuộc loại hình \textbf{bản dịch âm 1--1} (phiên âm từng chữ), không chứa các phân đoạn dịch nghĩa tự do. Về nguyên tắc, mỗi tự hình Hán hoặc Nôm trong câu đối ứng trực tiếp với một âm tiết Quốc ngữ theo đúng trật tự cú pháp tiếng Việt, nhờ đó cho phép định lượng mức độ tương đồng giữa hai vế dựa trên đối chiếu âm đọc. Tuy nhiên, quá trình chuyển hóa từ các trang ảnh quét lịch sử thành tập dữ liệu song song được dóng hàng chuẩn xác ở mức câu và mức tự hình đối mặt với nhiều khó khăn đan xen về mặt vật lý, thị giác và ngôn ngữ học; thay vì chỉ tối ưu riêng lẻ mô hình nhận dạng hay thuật toán so khớp chuỗi, bài toán đòi hỏi phải xử lý đồng bộ ngay từ khâu tái lập cấu trúc quyển sách cho tới bước chuẩn hóa mã ký tự và tri thức từ vựng chuyên biệt.

\subsection{Bố cục đối xứng ngược chiều và các dị thường vật lý của bản quét}

Về mặt trình bày, cả hai tập sách tuân theo quy ước đóng quyển song ngữ truyền thống tại Việt Nam đầu thế kỷ XX. Theo quy ước này, phần Quốc ngữ được bố trí ở nửa đầu cuốn sách, in ngang theo chữ Latinh từ trái sang phải, từ trên xuống dưới và đánh số trang tăng dần theo chiều xuôi. Ngược lại, phần Hán Nôm nằm ở nửa sau cuốn sách, khắc in theo cột dọc từ phải sang trái, trong mỗi cột đọc từ trên xuống dưới, đồng thời thứ tự trang đi ngược từ cuối sách lùi dần về giữa sách. Trong trường hợp lý tưởng khi tổng số trang chẵn đối xứng là $K$ và hai nửa không bị xáo trộn, trang Quốc ngữ thứ $p_{\text{vi}}$ sẽ ứng với trang Hán Nôm thứ $p_{\text{ch}} = K - p_{\text{vi}}$.

Tuy vậy, khi rà soát thực tế toàn bộ $1.321$ trang ảnh của hai tập, nhóm nhận thấy giả định tuyến tính đối xứng toàn cục bị gián đoạn bởi bốn dạng sai lệch vật lý cần được hiệu chỉnh trước khi đưa vào nhận dạng và ghi nhận lại trong bản PDF cập nhật của mỗi tập. Thứ nhất, ở Tập 2 xuất hiện hiện tượng quét đúp tại hai trang gốc thứ $27$ và $28$, vốn lặp lại nguyên văn nội dung của hai trang $25$ và $26$ liền trước; nếu không loại bỏ hai trang thừa này ngay từ đầu, chỉ số của toàn bộ phần còn lại của Tập 2 sẽ bị lệch đi hai đơn vị so với $683$ trang hiệu lực thực tế. Thứ hai, có bốn trang Hán Nôm bị đặt ngược $180^\circ$ khi số hóa, cụ thể là trang $421$ ở Tập 1 cùng ba trang $453, 485, 521$ ở Tập 2 (tính theo đánh số hiệu lực), làm cho số trang in nằm ngược ở cạnh dưới và đảo lộn toàn bộ chiều đọc cột. Thứ ba, ở Tập 1 xảy ra tình trạng hoán vị hai cặp tay sách khi đóng gáy, khiến hai trang Hán Nôm tương ứng với trang Quốc ngữ $43, 44$ bị xếp xuống vị trí $582, 581$, còn hai trang tương ứng với trang Quốc ngữ $53, 54$ lại nằm ở vị trí $592, 591$. Thứ tư, cũng trong Tập 1, hai trang Quốc ngữ liên tiếp mang số trang in sẵn $- 276 -$ và $- 277 -$ bị bỏ sót trong quá trình quét trong khi hai trang Hán Nôm đối ứng ($359$ và $358$) vẫn còn nguyên vẹn; sự khuyết thiếu đơn phía này tạo ra một bước nhảy pha $\Delta = 2$ làm thay đổi quy luật ghép cặp cho toàn bộ $39$ cặp trang cuối của Tập 1.

\subsection{Đặc trưng tự hình, ngữ âm và chính tả lịch sử}

Bên cạnh các sai lệch vật lý của bản quét, bản thân văn bản Công giáo Hán Nôm -- Quốc ngữ đầu thế kỷ XX cũng mang nhiều đặc trưng khiến các công cụ xử lý tiếng Việt hiện đại khó đạt độ chính xác cao nếu áp dụng trực tiếp. Trước hết, văn bản Hán Nôm không có khoảng trắng phân tách từ và cách dùng dấu khuyên tròn (\nom{。}) thường không tương ứng 1--1 với dấu chấm câu của vế Quốc ngữ; một câu Quốc ngữ có thể bao gồm nhiều vế ngắn ngăn bởi dấu khuyên tròn, hoặc ngược lại, nhiều đoạn Hán Nôm dài không hề có dấu ngắt câu. Nếu tách câu độc lập trên hai vế rồi mới ghép cặp, chỉ một vài điểm ngắt sai ở đầu trang cũng dễ dẫn tới trượt lệch câu kéo dài.

Thêm vào đó, hệ thống chữ viết trong bản khắc sử dụng nhiều chữ Nôm hiếm thuộc các mặt phẳng mở rộng của chuẩn Unicode (Extension B--I) và cả một số tự hình chưa được mã hóa chính thức. Do giới hạn của bảng mã cũ, công cụ nhận dạng chữ Nôm thường trả về các điểm mã thuộc vùng sử dụng riêng (PUA) không tra cứu được trong từ điển chuẩn, hoặc dễ nhận nhầm chữ Nôm có nét mờ sang chữ Hán phổ thông gần hình dáng. Song song với vấn đề mã hóa tự hình là lớp từ vựng tôn giáo mang quy ước ký âm đặc thù của văn hiến Công giáo thế kỷ XVII--XX: nhiều thuật ngữ thần học, tên thánh và địa danh gốc Latinh -- Bồ Đào Nha được ghi bằng chữ Hán Nôm mượn âm riêng biệt, chẳng hạn dùng \nom{秋} cho âm \textit{su} trong \textit{Giê-su}, hay dùng \nom{計} và \nom{移} cho các âm \textit{ghe}, \textit{rê} trong \textit{I-ghe-rê-gia} (\textit{Ecclesia}), vốn ít khi xuất hiện trong các từ điển Hán-Việt hay từ điển Nôm thông dụng. Cuối cùng, về phía bản in Quốc ngữ đầu thế kỷ XX, văn bản vẫn bảo lưu nhiều hình thái chính tả cổ như \textit{lở lại}, \textit{lọn lành}, \textit{vì chưng}, \textit{dường ấy} cùng dấu gạch nối nội tại trong tên riêng phiên âm. Trên nền giấy ố vàng với mực in có chỗ nhòe hoặc đứt nét, các tổ hợp hai tầng dấu tiếng Việt và nét gạch ngang mảnh của chữ \textit{đ} rất dễ bị nhầm sang \textit{d} hoặc bị các mô hình ngôn ngữ nhỏ tự ý sửa thành từ hiện đại.
